% !TeX root = ../main.tex
\chapter{Discussion générale}
\label{chap:analyse}

Ce chapitre synthétise l'ensemble des résultats obtenus, applique le critère de compatibilité métrologique entre les deux capteurs indépendants et analyse les principales sources d'erreurs expérimentales.

\section{Récapitulatif des mesures}
\label{sec:recap_resultats}

Le tableau~\ref{tab:recap-labo} rassemble les grandeurs physiques mesurées avec leurs incertitudes estimées.

\begin{table}[H]
  \centering
  \footnotesize
  \setlength{\tabcolsep}{3pt}
  \renewcommand{\arraystretch}{1.1}
  \begin{tabularx}{\textwidth}{l >{\raggedright\arraybackslash}X c c c}
    \toprule
    \textbf{Dispositif} & \textbf{Grandeur} & \textbf{Résultat mesuré} & \textbf{Référence / Condition} & \textbf{Incertitude / Écart} \\
    \midrule
    \multirow{2}{*}{CTN SEMI833} & Résistance ambiante $R_{\text{amb}}$ & \qty{97.20}{\kilo\ohm} & \qty{103.60}{\kilo\ohm} (\qty{20}{\degreeCelsius}) & Pente constructeur \\
     & Température ambiante $T_{\text{amb, NTC}}$ & \qty{21.54 \pm 0.80}{\degreeCelsius} & -- & Tolérance $\pm \qty{0.8}{\degreeCelsius}$ \\
    \midrule
    \multirow{5}{*}{Pt100 KP (chemisée)} & Résistance 4W ambiante & \qty{108.530}{\ohm} & \qty{100.00}{\ohm} (\qty{0}{\degreeCelsius}) & Multimètre 6½ digits \\
     & Température ambiante $T_{\text{amb, Pt100}}$ & \qty{22.16 \pm 0.15}{\degreeCelsius} & -- & Tolérance classe A ($\pm \qty{0.15}{\degreeCelsius}$) \\
     & Température LabVIEW & \qty{22.08 \pm 0.15}{\degreeCelsius} & \qty{22.16}{\degreeCelsius} & Écart : \qty{-0.08}{\degreeCelsius} \\
     & Résistance étalonnée $R_0$ (\qty{0}{\degreeCelsius}) & \qty{100.10 \pm 0.03}{\ohm} & \qty{100.00}{\ohm} & Écart : \qty{+0.10}{\ohm} (\qty{+0.1}{\percent}) \\
     & Biais de ligne 2W ($2 R_L$) & \qty{0.09}{\ohm} & \qty{0.00}{\ohm} (4W) & $\Delta T = \qty{+0.23}{\degreeCelsius}$ \\
    \midrule
    \multirow{2}{*}{Dynamique Pt100} & Constante de temps $\tau_{\text{fit}}$ & \qty{5.80 \pm 0.05}{\second} & \qty{5.80}{\second} (tangente) & $\text{RMSE} = \qty{0.09}{\degreeCelsius}$ \\
     & Temps de réponse à \qty{95}{\percent} & \qty{17.37 \pm 0.15}{\second} & $3\tau = \qty{17.40}{\second}$ & Écart : \qty{-0.2}{\percent} \\
    \midrule
    \multirow{2}{*}{Flux d'air CTN} & Palier standard ($N=497$) & \qty{41.13 \pm 0.13}{\degreeCelsius} & -- & Bruit : $\sigma = \qty{2.98}{\degreeCelsius}$ \\
     & Palier suréchantillonné ($N=352$) & \qty{42.83 \pm 0.17}{\degreeCelsius} & -- & Bruit : $\sigma = \qty{3.15}{\degreeCelsius}$ \\
    \midrule
    Auto-échauffement & Puissance $P$ / $\Delta T$ & \qty{0.061}{\milli\watt} & $\Delta T \approx \qty{0.09}{\degreeCelsius}$ & Négligeable ($< \qty{0.8}{\degreeCelsius}$) \\
    \bottomrule
  \end{tabularx}
  \caption{Synthèse des résultats expérimentaux et incertitudes associées.}
  \label{tab:recap-labo}
\end{table}

\section{Compatibilité entre la mesure CTN et la mesure Pt100}
\label{sec:compatibilite}

Pour vérifier la cohérence des deux mesures indépendantes de la température ambiante du laboratoire :
\begin{itemize}[leftmargin=1.2cm]
    \item CTN : $T_{\text{amb, NTC}} = \qty{21.54 \pm 0.80}{\degreeCelsius}$ (incertitude dominée par la tolérance constructeur de $\pm \qty{0.8}{\degreeCelsius}$) ;
    \item Pt100 : $T_{\text{amb, Pt100}} = \qty{22.16 \pm 0.15}{\degreeCelsius}$ (incertitude dominée par la classe A de $\pm \qty{0.15}{\degreeCelsius}$).
\end{itemize}
L'écart mesuré entre les deux capteurs est de :
\begin{equation}
    \Delta T = |21.54 - 22.16| = \qty{0.62}{\degreeCelsius}
\end{equation}
L'incertitude combinée vaut :
\begin{equation}
    u_c = \sqrt{u_{\text{NTC}}^2 + u_{\text{Pt100}}^2} = \sqrt{0.80^2 + 0.15^2} \approx \qty{0.814}{\degreeCelsius}
\end{equation}
Le facteur de compatibilité normalisé est :
\begin{equation}
    C = \frac{\Delta T}{u_c} = \frac{0.62}{0.814} \approx \mathbf{0.76}
\end{equation}
Comme $C \le 2$, l'écart de \qty{0.62}{\degreeCelsius} est statistiquement justifié par les incertitudes de mesure. Les deux méthodes sont en parfait accord métrologique.

\section{Analyse des principales sources d'erreurs}
\label{sec:sources_erreurs}

L'analyse met en évidence trois sources d'erreurs majeures :
\begin{enumerate}[leftmargin=1.2cm]
    \item \textbf{Résistance de ligne en 2 fils sur la Pt100 ($+\qty{0.23}{\degreeCelsius}$) :} L'étalonnage à \qty{0}{\degreeCelsius} a montré que le câble de liaison de \qty{1}{\meter} ajoute une résistance de $2 R_{\text{ligne}} = 100.16 - 100.07 = \qty{0.09}{\ohm}$. Avec une sensibilité de $\qty{0.385}{\ohm\per\degreeCelsius}$, cela provoque une erreur systématique positive de $\Delta T = \qty{+0.23}{\degreeCelsius}$, supérieure à la tolérance de classe A ($\pm \qty{0.15}{\degreeCelsius}$). Le montage 4 fils Kelvin est donc indispensable dès que la sonde est déportée.
    \item \textbf{Auto-échauffement de la CTN ($\approx \qty{0.09}{\degreeCelsius}$) :} Sous la tension de pont de \qty{14}{\volt} avec $R = \qty{462.35}{\kilo\ohm}$, la puissance dissipée dans la thermistance est d'environ \qty{0.061}{\milli\watt}. Avec un coefficient de dissipation thermique en air calme de $\delta \approx \qty{0.7}{\milli\watt\per\degreeCelsius}$, l'élévation interne de température est limitée à \qty{0.09}{\degreeCelsius}. Cette erreur représente environ \qty{11}{\percent} de la tolérance du capteur et valide le choix de la résistance série de \qty{470}{\kilo\ohm}.
    \item \textbf{Turbulence du flux d'air et non-répétabilité manuelle au sèche-cheveux :} 
    Le suréchantillonnage matériel (moyenne de 100 points) élimine efficacement le bruit haute fréquence de numérisation de la carte DAQ. Cependant, l'écart-type résiduel reste élevé ($\sigma \approx \qty{3}{\degreeCelsius}$) sur les deux séries. Cela montre que le bruit provient directement des tourbillons turbulents du jet d'air chaud et non de l'électronique.
    De plus, le décalage de température moyenne entre les deux essais ($\approx \qty{1.7}{\degreeCelsius}$) s'explique par la tenue manuelle : le sèche-cheveux et la perle n'étaient pas fixés sur un banc rigide, ce qui entraîne de légères variations de position et de distance entre les mesures.
\end{enumerate}
